% !TeX root = ../../Book1.tex
% 纲要标题：### 21.4 尺规作图
% 纲要位置：计划纲要.md，第 724 行。

\section{尺规作图}
\label{b1:sec:2104}

尺规作图限制了允许执行的操作：直尺只能连接已知点，圆规只能以已知点为圆心、以已知线段为半径画圆。把点写成坐标以后，直线与圆的交点只要求解一次或二次方程。因此，这种几何限制可以转化为扩张次数的限制。

% 纲要内容（仅作写作计划，尚非教材正文）：
%
% * 二次扩张塔与规矩数
% * 倍立方体、一般三等分角与正多边形
% * 化圆为方与 π 的超越性；超越性论证给出证明概要
%
% ---
%

\subsection{二次扩张塔与规矩数}
\label{b1:subsec:2104:01}

\begin{definition}\label{b1:def:constructible-number}
从平面上的 \((0,0)\)、\((1,0)\) 出发，允许有限次画过两个不同已知点的直线、以已知点为圆心且以两个已知点间距离为半径的圆，并把任意两条所画曲线的孤立交点加入已知点。重合曲线不由此指定新的点。如果一个点能经过有限次这种操作得到，则称它为\term[keguitu-dian]{尺规可作图点}[constructible point]。给定实数 \(a\)，如果 \((a,0)\) 是尺规可作图点，则称 \(a\) 为\term[guiju-shu]{规矩数}[constructible number]，也称可构造数、可作图数。\footnote{北京交通大学《近世代数》课程使用“规矩数”\cite[第 23 章第 1 节“几何作图简史、规矩数”]{BJTU2026ModernAlgebraCourse}；“可构造数”见席南华《基础代数》第一卷\cite[p. 125]{Xi2016BasicAlgebraI}，该书从复平面讨论可构造点。}\termidx[kegouzao-shu]{可构造数}\termidx[kezuotu-shu]{可作图数}
\end{definition}

标准的作垂线、平行线与转移线段长度均由上述操作实现。因此一个点可作图当且仅当它的两个坐标都是规矩数：从点向坐标轴作垂线得到坐标，反过来在两轴上标出坐标再作平行线，交点即为所需点。

\begin{theorem}\label{b1:thm:constructible-quadratic-tower}
实数 \(a\) 可作图，当且仅当存在全部位于 \(\mathbb R\) 中的域塔
\[
\mathbb Q=K_0\subseteq K_1\subseteq\cdots\subseteq K_r,
\qquad [K_i:K_{i-1}]=2,
\]
使 \(a\in K_r\)。允许 \(r=0\)。特别地，规矩数在 \(\mathbb Q\) 上的次数必为 \(2\) 的幂。
\end{theorem}
\begin{proof}
若当前全部已知坐标在实域 \(F\) 中，过已知点的直线方程系数属于 \(F\)。以已知点为圆心、已知线段为半径的圆方程是
\[
(x-u)^2+(y-v)^2=r^2,
\qquad u,v,r^2\in F.
\]
两直线的孤立交点由一次方程求出，坐标仍在 \(F\)。一条直线与一个圆联立，消去一个变量便得到至多二次方程；求出该变量后，另一个变量是它的 \(F\) 线性表达式。两圆相减得到一次方程，故也归结为前一种情况。重合的曲线不指定新的孤立交点；无交点时没有可执行步骤。因此每次新添交点至多需要一个实二次扩张。沿有限作图过程依次加入坐标，删去次数为 \(1\) 的步骤，就得到所需塔。

反过来，可作图的有向线段长度对加、减封闭，乘、除由相似三角形实现。例如连结 \((1,0)\) 与 \((0,b)\)，过 \((a,0)\) 作平行线，其纵轴截距为 \(ab\)。若 \(a\neq0\)，改为连结 \((a,0)\) 与 \((0,b)\)，过 \((1,0)\) 作平行线，纵轴截距便是 \(b/a\)。这些操作连同符号方向使规矩数构成一个实子域。

若 \(u>0\) 已可作图，在横轴上取 \(A=(0,0)\)、\(C=(1,0)\)、\(B=(1+u,0)\)，画以 \(AB\) 为直径的半圆，并过 \(C\) 作垂线，与上半圆交于 \(P=(1,h)\)，如\cref{b1:fig:square-root-semicircle}所示。圆心为 \(AB\) 的中点，半径为 \(AB\) 的一半，故
\[
\left(1-\frac{1+u}{2}\right)^2+h^2
=\left(\frac{1+u}{2}\right)^2,
\qquad h^2=u.
\]
取上半圆使 \(h>0\)，便作出了 \(\sqrt u\)。

% !TeX root = ../Book1.tex
\begin{figure}[htbp]
\centering
\begin{tikzpicture}[scale=1.35]
  % The proportions illustrate u=9/4; the construction works for every u>0.
  \pgfmathsetmacro{\diagramu}{2.25}
  \pgfmathsetmacro{\diagramr}{(1+\diagramu)/2}
  \pgfmathsetmacro{\diagramh}{sqrt(\diagramu)}
  \coordinate (A) at (0,0);
  \coordinate (C) at (1,0);
  \coordinate (B) at ({1+\diagramu},0);
  \coordinate (O) at (\diagramr,0);
  \coordinate (P) at (1,\diagramh);
  \draw[thick] (A)--(B);
  \draw[thick,color=AlgebraBlue] (A) arc[start angle=180,end angle=0,radius=\diagramr];
  \draw[thick,color=AlgebraBlue] (C)--(P);
  \draw[gray,dashed] (O)--(P);
  \draw (1,0.14)--(1.14,0.14)--(1.14,0);
  \foreach \point in {A,B,C,O,P} {\fill (\point) circle (1.2pt);}
  \node[below left] at (A) {\(A\)};
  \node[below] at (C) {\(C\)};
  \node[below right] at (B) {\(B\)};
  \node[below right] at (O) {\(O\)};
  \node[above left] at (P) {\(P\)};
  \node[below] at (0.5,-0.29) {\(1\)};
  \node[below] at ({1+\diagramu/2},-0.29) {\(u\)};
  \node[left,color=AlgebraBlue] at (0.92,{\diagramh/2}) {\(h\)};
\end{tikzpicture}
\caption[半圆开平方作图]{半圆开平方作图。\(AC=1\)、\(CB=u>0\)，\(O\) 为 \(AB\) 的中点，过 \(C\) 的垂线与上半圆交于 \(P\)。高度 \(CP=h\) 满足 \(h^2=u\)。图形比例仅作示意。}
\label{b1:fig:square-root-semicircle}
\end{figure}


最后，对任意实二次扩张 \(F(\beta)/F\)，把 \(\beta\) 的首一最小多项式写成 \(X^2+bX+c\)，配方得到 \((2\beta+b)^2=b^2-4c\)。判别式属于 \(F\) 且为正，所以 \(\beta=(-b\pm\sqrt{b^2-4c})/2\) 可由域运算和开正平方根得到。对塔归纳，末端全部元素都可作图。由\cref{b1:thm:tower-law}，\([K_r:\mathbb Q]=2^r\)，而 \(\bigl[\mathbb Q(a):\mathbb Q\bigr]\) 整除它，故为 \(2\) 的幂。
\end{proof}

若从更多给定点出发，应把 \(\mathbb Q\) 换为这些点坐标生成的域，证明逐字适用。次数为 \(2\) 的幂只是一个必要条件；完整判据要求存在二次扩张塔，不能省掉“塔”而只检查单个次数。

\subsection{倍立方体、三等分角与正多边形}
\label{b1:subsec:2104:02}

倍立方体要求从单位棱长作出棱长 \(\sqrt[3]{2}\)。\cref{b1:thm:eisenstein}对素元 \(2\) 给出 \(X^3-2\) 在 \(\mathbb Q\) 上不可约，故该数次数为 \(3\)，由\cref{b1:thm:constructible-quadratic-tower}不可作图。

\ProseParagraphBreak
一般三等分角也不可能，因为已经可作出的 \(60^\circ\) 角不能被尺规三等分。若 \(20^\circ\) 可作图，则 \(x=2\cos20^\circ\) 可作图；三倍角恒等式给出 \(x^3-3x-1=0\)。由\cref{b1:prop:rational-root}，这个首一整系数三次多项式的有理根只能是 \(\pm1\)，代入都非零。三次多项式可约必有一次因子，因此它不可约，\(x\) 的次数为 \(3\)，矛盾。这里排除的是一种对任意角都适用的作图方法；例如 \(90^\circ\) 可以三等分，并不与结论冲突。

\ProseParagraphBreak
正多边形问题的完整回答是\term[Gauss-Wantzel-dingli]{Gauss--Wantzel 定理}[Gauss--Wantzel theorem]。\footnote{旺策尔（Pierre Laurent Wantzel，1814-1848），法国数学家，证明若干古典尺规作图问题不可能完成。}

\begin{theorem}[Gauss--Wantzel]\label{b1:thm:constructible-polygon}
设 \(n\geq3\) 为整数。正 \(n\) 边形可尺规作图，当且仅当 \(\varphi(n)\) 为 \(2\) 的幂；等价地，
\[
n=2^a p_1\cdots p_s,
\]
其中 \(a,s\geq0\) 为整数，各 \(p_i\) 是两两不同的奇素数，并且每个 \(p_i\) 都可写成 \(2^{2^{r_i}}+1\)，其中 \(r_i\geq0\) 为整数；这样的素数称为\term[Fermat-sushu]{Fermat 素数}[Fermat prime]。\footnote{费马（Pierre de Fermat，1601？-1665），法国数学家，研究数论与解析几何；其出生年份有不同记载。}
\end{theorem}
\begin{proof}
正 \(n\) 边形可作图等价于 \(\zeta_n\) 的实部和虚部可作图。若可作图，把这两个坐标所在的实二次塔合并，再加入 \(i\)，得到一个次数为 \(2\) 的幂的域，其中含有 \(\zeta_n\)。由\cref{b1:thm:cyclotomic-irreducible}及塔公式，\(\varphi(n)=\bigl[\mathbb Q(\zeta_n):\mathbb Q\bigr]\) 也是 \(2\) 的幂。

反过来，若 \(\varphi(n)\) 为 \(2\) 的幂，\cref{b1:thm:cyclotomic-galois-group}说明分圆域的 Galois 群是有限交换 \(2\) 群。复共轭子群的固定域是实子域 \(F=\mathbb Q(\zeta_n+\zeta_n^{-1})\)：这里 \(n\geq3\)，故 \(\zeta_n\) 不实，且它在 \(F\) 上满足 \(X^2-(\zeta_n+\zeta_n^{-1})X+1=0\)。复共轭子群在交换群中正规，所以\cref{b1:thm:galois-quotient}说明 \(F/\mathbb Q\) 也是 Galois 扩张，其群是有限交换 \(2\) 群。这样的群有一条各步指数均为 \(2\) 的子群列：对非平凡有限交换 \(2\) 群取极大真子群，商为交换单群，由\cref{b1:prop:abelian-simple}为 \(C_2\)，再对阶较小的子群归纳。由\cref{b1:thm:finite-galois-correspondence}，\(F\) 因而位于一条实二次扩张塔的末端。于是 \(\cos(2\pi/n)\) 可作图，\(\sin(2\pi/n)\) 也可由 \(\sqrt{1-\cos^2(2\pi/n)}\) 作出，正 \(n\) 边形便可作图。

最后设 \(n=2^a\prod p^{e_p}\)，其中 \(p\) 为奇素数。由模素数幂计数及中国剩余定理，
\[
\varphi(n)=\varphi(2^a)\prod_{p\mid n,\ p\text{ 奇}}p^{e_p-1}(p-1).
\]
它为 \(2\) 的幂当且仅当每个奇素数的指数 \(e_p=1\)，且 \(p-1\) 是 \(2\) 的幂。若 \(p=2^m+1\) 是素数，写 \(m=2^r u\)，其中 \(u\) 为奇数。若 \(u>1\)，则 \(z+1\) 整除 \(z^u+1\)，取 \(z=2^{2^r}\) 会给出 \(p\) 的非平凡因子。因此 \(u=1\)，即 \(m=2^r\)。反向代入上式立即得到 \(\varphi(n)\) 为 \(2\) 的幂。
\end{proof}

例如正 \(3,5,17\) 边形及正 \(15\) 边形可作图，正 \(7\) 边形与正 \(9\) 边形不可作图，因为相应次数为 \(6\)。每个给定的 \(n\) 都可以按其素因数分解单独检查。

\subsection{化圆为方与圆周率的超越性}
\label{b1:subsec:2104:03}

化圆为方要求作出与单位圆等面积的正方形，其边长必须为 \(\sqrt\pi\)。为排除这种长度，需要研究指数函数在代数数处的取值。下述结果是\term[Lindemann-Weierstrass-dingli]{Lindemann--Weierstrass 定理}[Lindemann--Weierstrass theorem]的单变量推论\footnote{林德曼（Carl Louis Ferdinand von Lindemann，1852-1939），德国数学家，证明圆周率的超越性，从而解决化圆为方的不可能性。}；其证明方法将整数的离散性与指数积分的估计结合起来\footnote{魏尔斯特拉斯（Karl Theodor Wilhelm Weierstrass，1815-1897），德国数学家，发展复函数论与分析的严格基础。}。

\begin{theorem}[Lindemann--Weierstrass 的单变量推论]\label{b1:thm:lindemann-weierstrass-special-case}
若 \(\alpha\) 是非零代数数，则 \(e^\alpha\) 是超越数。
\end{theorem}
\begin{proofsketch}
假设 \(e^\alpha=\beta\) 为代数数。取非零多项式 \(P(T)\in\mathbb Z[T]\)，使 \(P(\beta)=0\) 且 \(P(0)\ne0\)，并记 \(\alpha_1=\alpha,\ldots,\alpha_r\) 为 \(\alpha\) 的全部共轭。令
\[
\Lambda=\mathbb Z\alpha_1+\cdots+\mathbb Z\alpha_r\subseteq(\mathbb C,+).
\]
这是有限生成无挠交换群，\cref{b1:thm:fga-structure}保证它有一组整数基 \(\lambda_1,\ldots,\lambda_s\)：每个 \(\gamma\in\Lambda\) 唯一写成 \(\sum_j a_j\lambda_j\)，其中 \(a_j\in\mathbb Z\)。其整群代数 \(\mathbb Z[\Lambda]\) 由形式有限和 \(\sum_\gamma b_\gamma E^\gamma\) 组成，乘法规定为 \(E^\gamma E^\eta=E^{\gamma+\eta}\)。整数基给出同构
\[
\mathbb Z[\Lambda]\cong\mathbb Z[T_1^{\pm1},\ldots,T_s^{\pm1}],
\qquad E^{\sum_j a_j\lambda_j}\longmapsto\prod_j T_j^{a_j}.
\]
右端是 Laurent 多项式环，因而是整环。在这个形式环中构造
\[
Q=\prod_{i=1}^r P(E^{\alpha_i}).
\]
每个 \(\alpha_i\ne0\)，故 \(P(E^{\alpha_i})\) 中不同次数对应不同形式指数，是非零元素；整环性质于是给出 \(Q\ne0\)。可写 \(Q=\sum_\gamma b_\gamma E^\gamma\)，其中 \(b_\gamma\in\mathbb Z\)。代数共轭置换保留 \(Q\)。现在使用求值同态
\[
\operatorname{ev}:\mathbb Z[\Lambda]\longrightarrow\mathbb C,
\qquad E^\gamma\longmapsto e^\gamma,
\]
在 Laurent 多项式坐标中即 \(T_j\mapsto e^{\lambda_j}\)。由 \(P(e^\alpha)=0\)，\(Q\) 的求值为零。这一步并未假定求值同态单射：形式环中的非零元素是否可能求值为零，正是这里要排除的情形。

在形式环中令 \(Q^-=\sum_\gamma b_\gamma E^{-\gamma}\)，即把 \(Q\) 的各指数取负，系数不变。上标 \(-\) 表示这个操作，既不表示乘法逆，也不表示复共轭。乘积 \(QQ^-\) 的常数项为 \(\sum_\gamma b_\gamma^2>0\)，因为一项的指数与另一项的负指数相消，恰好要求原来的两个指数相等。\(QQ^-\) 仍在代数共轭下保持，求值仍为零。合并相同指数，得到
\[
c_0+\sum_{j=1}^n c_j e^{z_j}=0,
\]
其中 \(c_j\) 为整数，\(c_0>0\)，各 \(z_j\) 为互异的非零代数数；指数集合在共轭下保持，且共轭置换保留相应系数。这里共轭只作用于形式表达式的指数；取通常的指数值是随后单独进行的代入。

余下的目标是构造一族不全为零的代数数 \(S_{m,i}\)，使它们乘以一个可控制的整数因子后成为代数整数，同时使全部代数共轭都足够小。非零代数整数的范数是非零整数，绝对值至少为 \(1\)；若所有共轭的乘积反而小于 \(1\)，便得到矛盾。这里约束的是范数，不是某一个共轭的绝对值。

令 \(z_0=0\)、\(f(z)=\prod_{j=0}^n(z-z_j)\)。在各 \(z_j\) 处安排高重零点，可以使许多低阶导数消失；除以 \((m-1)!\) 后，剩余导数中的阶乘比仍为整数，便于控制整性。为把这些导数与一个小积分联系起来，取全部导数之和，使相邻导数相减时消项。具体地，对正整数 \(m\) 及 \(0\le i\le n\)，记
\[
P_{m,i}(z)=\frac{f(z)^m}{z-z_i},\qquad
d_m=(n+1)m-1,\qquad
H_{m,i}(z)=\frac1{(m-1)!}\sum_{r=0}^{d_m}P_{m,i}^{(r)}(z).
\]
这里 \(P_{m,i}\) 是次数为 \(d_m\) 的多项式；它在 \(z_i\) 处有 \(m-1\) 重零，在其余 \(z_k\) 处有 \(m\) 重零。除去一个因子使这些多项式的次数小于 \((n+1)m\)，同时保留后面证明非奇异性所需的重零点。定义端点矩阵 \(A_m=(H_{m,i}(z_k))_{0\le i,k\le n}\)，并沿从 \(0\) 到 \(z_k\) 的线段取积分
\[
\varepsilon_{m,i}(z_k)
=\frac1{(m-1)!}\int_0^{z_k}e^{-z}P_{m,i}(z)\,\mathrm{d}z.
\]
这里的线段积分可用实参数直接表示：将 \(z=tz_k\)、\(0\le t\le1\) 代入，有
\[
\int_0^{z_k}F(z)\,\mathrm{d}z
=z_k\int_0^1F(tz_k)\,\mathrm{d}t,
\]
右端按实部、虚部分别作普通实积分；\(z_k=0\) 时两端均为零。因此积分的绝对值至多为 \(\lvert z_k\rvert\max_{0\le t\le1}\lvert F(tz_k)\rvert\)。
由于 \(H_{m,i}-H_{m,i}'=P_{m,i}/(m-1)!\)，分部积分给出可逐项核对的端点等式
\[
\varepsilon_{m,i}(z_k)
=H_{m,i}(0)-e^{-z_k}H_{m,i}(z_k).
\]

矩阵 \(A_m\) 非奇异，这一点不需要预设其行列式的精确公式。设 \(\sum_i\lambda_iH_{m,i}(z_k)=0\) 对全部 \(k\) 成立，令 \(H=\sum_i\lambda_iH_{m,i}\) 及 \(G=H-H'\)。由定义，\(G\) 被 \(f^{m-1}\) 整除，故在每个 \(z_k\) 处至少有 \(m-1\) 重零。利用 \(H'=H-G\) 逐次求导，并从 \(H(z_k)=0\) 出发归纳，可知 \(H\) 在每个 \(z_k\) 处至少有 \(m\) 重零。可是 \(\deg H\le d_m<(n+1)m\)，所以 \(H=0\)。再由
\[
G(z)=\frac{f(z)^{m-1}}{(m-1)!}
\sum_{i=0}^n\lambda_i\prod_{j\ne i}(z-z_j)
\]
可将括号内的多项式依次在 \(z_i\) 处取值：由于各 \(z_j\) 不同，每次只有第 \(i\) 项不为零，故全部 \(\lambda_i=0\)。

令 \(c=(c_0,\ldots,c_n)^{\mathsf T}\)，并设 \(S_{m,i}=(A_mc)_i=\sum_k c_kH_{m,i}(z_k)\)。由于 \(c_0>0\) 且 \(A_m\) 非奇异，这些代数数不全为零。把端点等式乘以 \(c_ke^{z_k}\) 后求和，由 \(\sum_k c_ke^{z_k}=0\) 得
\[
S_{m,i}=-\sum_{k=0}^n c_ke^{z_k}\varepsilon_{m,i}(z_k).
\]
线段只有有限条；其上 \(e^{-z}\)、\(f(z)\) 与各多项式 \(f(z)/(z-z_i)\) 有共同的有限上界。写 \(P_{m,i}=f^{m-1}\bigl(f/(z-z_i)\bigr)\)，并使用上面的实参数积分界，便得到与 \(m,i\) 无关的 \(C,B>0\)，使 \(\lvert S_{m,i}\rvert\le CB^m/(m-1)!\)。

取正整数 \(q\)，使每个 \(qz_j\) 都是代数整数。\(P_{m,i}\) 在 \(z_k\) 处至少有 \(m-1\) 重零；展开 \(P_{m,i}(z_k+t)\) 后，\(t^r\) 的系数是含 \(d_m-r\) 个差 \(z_k-z_j\) 的整数系数多项式，乘以 \(q^{d_m-r}\) 即成为代数整数。对非零导数项，\(P_{m,i}^{(r)}(z_k)/(m-1)!\) 中的阶乘比 \(r!/(m-1)!\) 是整数。因此取固定的 \(D=q^{n+1}\)，则 \(D^mH_{m,i}(z_k)\)，从而 \(D^mS_{m,i}\)，均为代数整数。

取包含全部 \(z_j\) 的固定 Galois 数域 \(M\)。指数 \(z_j\) 及对应系数 \(c_j\) 在代数共轭下同步置换；若 \(\sigma(z_i)=z_j\)，则 \(\sigma\) 把 \(P_{m,i}\)、\(H_{m,i}\) 分别变为 \(P_{m,j}\)、\(H_{m,j}\)，从而 \(\sigma(S_{m,i})=S_{m,j}\)。故上面的积分界同时控制全部共轭。若 \(D^mS_{m,i}\ne0\)，其范数 \(N_{M/\mathbb Q}(D^mS_{m,i})\) 是非零整数；但其绝对值至多为 \((CD^mB^m/(m-1)!)^{[M:\mathbb Q]}\)，右边随 \(m\) 趋于零。故充分大的 \(m\) 下全部 \(S_{m,i}=0\)。这与它们不全为零矛盾。

原文\cite[\S 1, pp. 214--216]{Lindemann1882Pi}使用 Hermite 积分\footnote{埃尔米特（Charles Hermite，1822-1901），法国数学家，证明自然常数 \(e\) 的超越性，并研究椭圆函数。}及另行构造的系数矩阵；这里直接证明端点矩阵的非奇异性，不沿用原文系数矩阵的精确行列式等式。共轭对称与积分估计在原文\cite[\S\S 2--4, pp. 216--225]{Lindemann1882Pi}中用于推出单变量结论。这个证明概要使用复指数函数、线段上的积分估计和代数整数范数，不能仅由 Galois 对应替代。
\end{proofsketch}

一般定理断言，互异代数数的指数在代数数域上线性无关。取指数 \(0,\alpha\)，若 \(\beta=e^\alpha\) 为代数数，就有代数系数关系 \(-\beta e^0+e^\alpha=0\)，因此上述结论确是它的单变量推论。

\begin{corollary}\label{b1:cor:pi-transcendental}
圆周率 \(\pi\) 是超越数，因而不能用尺规将圆化为等面积的正方形。
\end{corollary}
\begin{proof}
若 \(\pi\) 代数，则非零数 \(i\pi\) 也代数，由\cref{b1:thm:lindemann-weierstrass-special-case}，\(e^{i\pi}\) 应当超越。这与 Euler 恒等式 \(e^{i\pi}=-1\) 矛盾，故 \(\pi\) 超越。

若 \(\sqrt\pi\) 是规矩数，由\cref{b1:thm:constructible-quadratic-tower}，它在 \(\mathbb Q\) 上代数；平方后 \(\pi\) 也应代数，仍然矛盾。因此化圆为方不可能。
\end{proof}
在化圆为方问题中，作图判据先给出可作图数必须代数这一必要条件；超越性论证另使用复指数函数及积分估计。Euler 恒等式也属于复指数函数的性质，因此化圆为方的不可作图性同时包含代数与分析的论证。
